% Bilaga C, satt ur project/README.md och info/examination.md.

\chapter{Grupprojektet}
\label{app:project}

Det här är projektets specifikation, samlad på ett ställe: arbetsformen, strukturen, modulernas
gränssnitt, testplanen och bedömningen. \Chapref{c10:ch} till \ref{c19:ch} är genomgången; det här är
kontraktet, och där de två skiljer sig gäller det som står här.

Uppgiften är att konstruera, simulera och köra en förenklad CAN-kontroller i VHDL, och göra den
nåbar från en mikrokontroller över SPI. Parallellklassen skriver C++-drivrutinen mot samma register,
så det som byggs är ena halvan av ett riktigt hårdvaru- och mjukvarusystem, inte en övning som slutar
i simulatorn.

När projektet är klart har gruppen implementerat en \textbf{bittimer} som producerar sampel- och
bitslutspulser, en \textbf{CRC-15-motor} som både beräknar och verifierar checksumman, ett
\textbf{sändningsskiftregister} som serialiserar och sätter in stoppbitar, ett
\textbf{mottagningsskiftregister} som deserialiserar och tar bort dem, en \textbf{CAN-kontroller}
vars tillståndsmaskin sekvenserar ramens fält och hanterar arbitrering, en \textbf{registerbank} som
översätter kontrollerns pulser till pollbara register, en \textbf{SPI-brygga} som gör registren
nåbara utifrån, och en \textbf{toppnivå} som binder ihop alltihop till en nod en mikrokontroller kan
prata med.

% ------------------------------------------------------------------------------------------
\section{Arbetsform}
\label{app:project:work}

Projektet genomförs i grupper om \textbf{4--5 studenter}. Gruppen ansvarar gemensamt för hela
projektet, och arbetssättet följer industriell praxis.

\subsection{Git och GitHub}
\begin{itemize}
  \item Skapa ett privat repository och bjud in alla gruppmedlemmar samt handledaren som
    collaborator.
  \item Arbeta i feature-branchar. En branch per modul är en rimlig utgångspunkt.
  \item Committa ofta, med meddelanden som säger vad ändringen gör och inte bara att den finns.
  \item \textbf{Pusha aldrig direkt till \code{main}.} All kod når \code{main} genom en Pull
    Request. Skydda branchen i GitHubs inställningar så att regeln inte beror på att alla kommer
    ihåg den.
\end{itemize}

\subsection{Kodgranskning}
\begin{itemize}
  \item Varje PR granskas av \textbf{minst en annan gruppmedlem} innan den mergas.
  \item Granskaren kontrollerar: att portordningen stämmer mot gränssnittstabellen, att modulens
    testbänk passerar, att inga värden är hårdkodade förbi \code{can_def}, och att varje entitets
    gränssnitt och varje viktig intern signal är kommenterad.
  \item Lägg granskningskommentarer i GitHub, inte i chatt eller muntligt. En granskning som inte
    finns skriftligt har inte hänt.
  \item En granskning som bara säger ”ser bra ut” är ingen granskning. Ställ minst en riktig fråga,
    även när koden är bra: \emph{varför} den gör som den gör är oftare intressant än \emph{att} den
    gör det.
\end{itemize}

\subsection{Arbetsfördelning}
Alla ska skriva VHDL och alla ska granska VHDL. Fördelningen får vara ojämn i mängd, inte i art.
Dokumentera fördelningen löpande i \file{CONTRIBUTORS.md} i repots rot.

% ------------------------------------------------------------------------------------------
\section{Så här drivs projektet}
\label{app:project:drive}

Projektet har \textbf{inga milstolpar}. Varje kapitel från \chapref{c10:ch} och framåt anger ett
\emph{riktvärde}, formulerat som ”ni bör vara ungefär här”. Riktvärdena är inte krav och betygsätts
inte.

Det är ett medvetet val och inte generositet. Efter \file{can_def.vhd} har de fyra lövmodulerna
\code{bit_timer}, \code{crc15}, \code{tx_shift_reg} och \code{rx_shift_reg} \textbf{inga beroenden
till varandra}, och var och en har sin egen utdelade testbänk. En grupp om 4--5 kan därför bygga dem
parallellt, i vilken ordning som helst. Att låsa fast en ordning skulle bara hindra det.

Men beroendegrafen är inte platt, och det är värt att planera efter:

\begin{textdrawing}
             meta_prev                                    (läser inget paket alls)
can_def  ──> bit_timer, crc15, tx_shift_reg, rx_shift_reg (de fyra löven, parallella)
                                    └──> can_controller   (kap 16 sändning, kap 17 mottagning)
can_def  ──> register_bank                                (oberoende, byggbar från kap 10)
spi_def  ──> spi_reg_bridge                               (oberoende, byggbar från kap 10)
spi_slave                                                 (utdelad, skrivs inte av gruppen)
    can_controller + register_bank + spi_reg_bridge + spi_slave ──> can_spi_node   (kap 19)
\end{textdrawing}

Två saker följer av den:

\begin{itemize}
  \item \textbf{\code{can_controller} är den kritiska vägen.} Den har ingen intern parallellitet,
    den sträcker sig över två pass, och dess testbänk är en systemtestbänk som antingen går igenom
    eller inte \textemdash{} ingen delvis grön. Den som ska äga den bör börja läsa \chapref{c16:ch}
    tidigt, och gruppen bör räkna med att den tar mest tid av allt.
  \item \textbf{\code{register_bank} och \code{spi_reg_bridge} är två fullvärdiga arbetsspår som kan
    starta direkt.} Den ena behöver bara \code{can_def}, den andra bara \code{spi_def}, och båda har
    utdelade testbänkar. Ingen behöver vänta på kontrollern för att ha något att göra, och en grupp
    som skjuter dem till \chapref{c18:ch} och \ref{c19:ch} har lagt sitt tightaste arbete sist.
\end{itemize}

Två hållpunkter är obligatoriska:

\begin{booktable}{lL}
  \thead{När} & \thead{Vad} \\
  \midrule
  \Chapref{c15:ch} & \textbf{Kodgranskningsseminarium.} Varje grupp visar en mergad PR och en
    granskning gruppen skrivit, och berättar om en sak i arbetssättet som inte fungerat. \\
  \Chapref{c19:ch} & \textbf{Slutredovisning och inlämning.} Gröna testbänkar, genomgång av
    \code{can_spi_node}, bring-up på hårdvara så långt gruppen kommit, och repots historik. \\
\end{booktable}

Det som bedöms är slutresultatet och arbetssättet: designen, Git-historiken, PR:erna och
granskningarna. \textbf{Inte takten.} En grupp som ligger efter halvvägs igenom projektet och lämnar
in en komplett, välgranskad nod har gjort projektet rätt.

% ------------------------------------------------------------------------------------------
\section{Projektstruktur}
\label{app:project:layout}

Två platta kataloger, med modulerna bredvid de testbänkar som kontrollerar dem:

\begin{textdrawing}
controller/   can_def.vhd, meta_prev.vhd, bit_timer.vhd, crc15.vhd,
              tx_shift_reg.vhd, rx_shift_reg.vhd, can_controller.vhd
              + de sex utdelade *_tb.vhd
bridge/       register_bank.vhd, spi_reg_bridge.vhd, can_spi_node.vhd
              + utdelade spi_slave.vhd, spi_def.vhd,
                register_bank_tb.vhd, spi_reg_bridge_tb.vhd
ci/           build_project.sh, kopierad från kursrepot
Makefile      med målet build-project
CONTRIBUTORS.md
.gitignore
\end{textdrawing}

Gränsen mellan de två katalogerna går vid protokollet: allt i \file{controller/} känner till CAN och
ingenting om SPI, och \textbf{ingen modul i någondera katalogen känner till båda}.
\code{register_bank} är mitten: den känner varken CAN-ramformatet eller SPI utan bara
registersemantik, och läser \code{can_def} enbart för porttyperna.

De utdelade filerna hämtas från kursrepot och läggs in oförändrade. \file{spi_slave.vhd} skriver
gruppen inte: den är transport, inte kursens ämne, och den läses i \chapref{c18:ch} i stället.

% ------------------------------------------------------------------------------------------
\section{Namnkonventioner och kodstil}
\label{app:project:style}

All kod och alla kommentarer i koden skrivs på \textbf{engelska}. Instruktionerna är på svenska.

\begin{itemize}
  \item Modul- och signalnamn i \code{snake_case}, till exempel \code{timer_enable}.
  \item Konstanter i versaler, till exempel \code{TICKS_PER_BIT}.
  \item \code{std_logic} och \code{std_logic_vector} för allt som är ledningar; \code{natural} för
    räknare.
  \item Inga hårdkodade värden. Allt som är protokoll ligger i \code{can_def}.
  \item Kommentera varje entitets gränssnitt och varje viktig intern signal.
  \item Håll gränssnitten stabila. Testbänkarna är utdelade och binder positionellt.
\end{itemize}

\begin{warning}
\textbf{Positionell portbindning.} Varje utdelad testbänk binder till sin modul \textbf{positionellt},
och det gör varje instansiering inuti \code{can_controller} också. Kontraktet är alltså
\textbf{ordningen och typerna} på portarna, exakt som tabellerna i \secref{app:project:controller}
och \ref{app:project:bridge} listar dem. Portarnas \emph{namn} är gruppens egna.

Byt plats på två portar av samma typ och ingenting klagar vid analys; konstruktionen beter sig bara
fel i simuleringen. Det är priset för att gränssnittstabellen får vara enda sanningskälla.
\end{warning}

% ------------------------------------------------------------------------------------------
\section{Det delade paketet \code{can_def}}
\label{app:project:candef}

Skrivs i \chapref{c10:ch}, i \file{controller/can_def.vhd}, och läses av allt annat.

\textbf{Konstanter:} \code{CLOCK_FREQ_HZ}, \code{BIT_RATE_HZ}, \code{TICKS_PER_BIT},
\code{SAMPLE_TICK}, \code{ID_WIDTH}, \code{DLC_WIDTH}, \code{DLC_MAX}, \code{CRC_WIDTH},
\code{DATA_WIDTH}, \code{MAX_RUN}, \code{CRC_POLY}.

\textbf{Subtyper:} \code{byte_t}, \code{id_t}, \code{crc_t}, \code{dlc_t}, \code{data_t}.

Riktvärden för ett DE0-CV med 50 MHz systemklocka och CAN på 1 Mbit/s:

\begin{booktable}{llL}
  \thead{Konstant} & \thead{Värde} & \thead{Kommentar} \\
  \midrule
  \code{CLOCK_FREQ_HZ} & \code{50_000_000} & Systemklockan. \\
  \code{BIT_RATE_HZ} & \code{1_000_000} & 1 Mbit/s. \\
  \code{TICKS_PER_BIT} & \code{CLOCK_FREQ_HZ / BIT_RATE_HZ} & 50 klockcykler per CAN-bit. Räkna ut
    den, skriv den inte. \\
  \code{SAMPLE_TICK} & \code{TICKS_PER_BIT * 7 / 10} & Sampelpunkten, ungefär 70 \% in i
    bitperioden. \\
  \code{ID_WIDTH} & \code{11} & Standardidentifierare. \\
  \code{DLC_WIDTH} & \code{4} & DLC är fyra bitar. \\
  \code{DLC_MAX} & \code{8} & Högst åtta databytes. \\
  \code{DATA_WIDTH} & \code{DLC_MAX * 8} & 64 bitar nyttolast. \\
  \code{CRC_WIDTH} & \code{15} & CAN:s CRC-15. \\
  \code{MAX_RUN} & \code{5} & Femregeln för bitstoppning. \\
  \code{CRC_POLY} & \code{"100010110011001"} & Motsvarar $x^{15} + x^{14} + x^{10} + x^{8} + x^{7} +
    x^{4} + x^{3} + 1$. \\
\end{booktable}

Att \code{TICKS_PER_BIT} och \code{SAMPLE_TICK} är \emph{uträknade} och inte inskrivna spelar roll:
byter någon bithastighet ska allt följa med.

% ------------------------------------------------------------------------------------------
\section{Modulerna i \code{controller/}}
\label{app:project:controller}

Kapitlen beskriver modulerna i detalj; det här är portordningen samlad, eftersom det är den som är
kontraktet.

\subsection*{\code{meta_prev} (\chapref{c12:ch})}
En generic, \code{WIDTH} (\code{natural := 1}), och tre portar: \code{clock} (in), \code{async_in}
(in, \code{std_logic_vector(WIDTH-1 downto 0)}), \code{sync_out} (out, samma typ).
Dubbelvippsynkroniseraren från \chapref{c4:ch}, skriven en gång med en breddgeneric. Den behöver
ingen reset: efter två klockcykler är innehållet ändå ersatt. Instansieras två gånger inuti
\code{can_controller}, för \code{reset_n} och för \code{rx_bus}.

\subsection*{\code{bit_timer} (\chapref{c12:ch})}
\begin{longbooktable}{cll>{\raggedright\arraybackslash}p{0.24\linewidth}>{\raggedright\arraybackslash}p{0.30\linewidth}}{\thead{\#} & \thead{Port} & \thead{Riktning} & \thead{Typ} &
  \thead{Betydelse}}
  1 & \code{clock} & in & \code{std_logic} & Systemklocka. \\
  2 & \code{reset_s2_n} & in & \code{std_logic} & Synkroniserad, aktiv låg reset. \\
  3 & \code{enable} & in & \code{std_logic} & Räknar medan \code{'1'}, håller medan \code{'0'}. \\
  4 & \code{resync} & in & \code{std_logic} & Tvingar tillbaka räknaren till tick 0. \\
  5 & \code{sample} & out & \code{std_logic} & Puls vid sampelpunkten. \\
  6 & \code{bit_done} & out & \code{std_logic} & Puls vid bitperiodens slut. \\
\end{longbooktable}

En räknare från \code{0} till \code{TICKS_PER_BIT-1} som pulsar \code{sample} vid \code{SAMPLE_TICK}
och \code{bit_done} vid periodens slut. \code{resync} används en gång per ram, vid SOF; däremellan
löper timern fritt.

\subsection*{\code{crc15} (\chapref{c13:ch})}
\begin{longbooktable}{cll>{\raggedright\arraybackslash}p{0.24\linewidth}>{\raggedright\arraybackslash}p{0.30\linewidth}}{\thead{\#} & \thead{Port} & \thead{Riktning} & \thead{Typ} &
  \thead{Betydelse}}
  1 & \code{clock} & in & \code{std_logic} & Systemklocka. \\
  2 & \code{reset_s2_n} & in & \code{std_logic} & Synkroniserad, aktiv låg reset. \\
  3 & \code{clear} & in & \code{std_logic} & Nollställer registret utan reset. \\
  4 & \code{enable} & in & \code{std_logic} & Integrerar \code{data} denna flank. \\
  5 & \code{data} & in & \code{std_logic} & Aktuell bit. \\
  6 & \code{crc} & out & \code{crc_t} & Registrets nuvarande värde. \\
  7 & \code{valid} & out & \code{std_logic} & \code{'1'} när registret är idel nollor. \\
\end{longbooktable}

Ett 15-bitars register med återkoppling. Samma motor både genererar och kontrollerar: sändaren matar
in sina bitar och läser ut \code{crc}; mottagaren matar in samma bitar \emph{plus} den mottagna
CRC:n, och \code{valid} går hög om allt stämmer. Ingen lägesomkoppling behövs. \code{clear} behövs
ändå: en avbruten ram lämnar skräp i registret.

\subsection*{\code{tx_shift_reg} (\chapref{c14:ch})}
\begin{longbooktable}{cll>{\raggedright\arraybackslash}p{0.24\linewidth}>{\raggedright\arraybackslash}p{0.30\linewidth}}{\thead{\#} & \thead{Port} & \thead{Riktning} & \thead{Typ} &
  \thead{Betydelse}}
  1 & \code{clock} & in & \code{std_logic} & Systemklocka. \\
  2 & \code{reset_s2_n} & in & \code{std_logic} & Synkroniserad, aktiv låg reset. \\
  3 & \code{load} & in & \code{std_logic} & Laddar \code{data}/\code{bit_count} \textbf{och} lägger
    ut första biten samma cykel. \\
  4 & \code{data} & in & \code{byte_t} & Gruppen, med sina giltiga bitar högst. \\
  5 & \code{bit_count} & in & \code{std_logic_vector(3 downto 0)} & Antal giltiga bitar, 1--8. \\
  6 & \code{shift} & in & \code{std_logic} & En puls per CAN-bit, från \code{bit_timer.bit_done}. \\
  7 & \code{tx_bit} & out & \code{std_logic} & Biten som ska drivas nästa bitperiod. \\
  8 & \code{stuff} & out & \code{std_logic} & \code{'1'} på just den skiftning som lägger ut en
    stoppbit. \\
  9 & \code{done} & out & \code{std_logic} & Pulsar när gruppen är helt utsänd. \\
  10 & \code{bit_valid} & out & \code{std_logic} & \code{'1'} den cykel en ny bit lades ut. \\
\end{longbooktable}

Skiftar ut MSB först, räknar lika bitar i rad, och sätter in en stoppbit med motsatt värde efter
\code{MAX_RUN} lika. \textbf{Stoppningstillståndet lever kvar över omladdningar}: regeln gäller
bitströmmen, inte den enskilda gruppen. \code{stuff} och \code{bit_valid} finns för
\code{can_controller}s skull: CRC:n matas på \code{bit_valid and not stuff}, eftersom en stoppbit
inte ingår i checksumman.

\subsection*{\code{rx_shift_reg} (\chapref{c15:ch})}
\begin{longbooktable}{cll>{\raggedright\arraybackslash}p{0.24\linewidth}>{\raggedright\arraybackslash}p{0.30\linewidth}}{\thead{\#} & \thead{Port} & \thead{Riktning} & \thead{Typ} &
  \thead{Betydelse}}
  1 & \code{clock} & in & \code{std_logic} & Systemklocka. \\
  2 & \code{reset_s2_n} & in & \code{std_logic} & Synkroniserad, aktiv låg reset. \\
  3 & \code{sample} & in & \code{std_logic} & En puls per CAN-bit, från \code{bit_timer.sample}. \\
  4 & \code{rx_bus} & in & \code{std_logic} & Bussnivån, giltig vid sampelpunkten. \\
  5 & \code{bit_count} & in & \code{std_logic_vector(3 downto 0)} & Antal \emph{riktiga} bitar att
    samla in, 1--8. \\
  6 & \code{enable} & in & \code{std_logic} & Avgör om \code{sample}-pulser behandlas. \\
  7 & \code{data} & out & \code{byte_t} & Den insamlade gruppen, högerjusterad. \\
  8 & \code{valid} & out & \code{std_logic} & Pulsar när \code{bit_count} riktiga bitar samlats
    in. \\
  9 & \code{stuff_error} & out & \code{std_logic} & Pulsar om en väntad stoppbit inte inverterar. \\
  10 & \code{real_bit} & out & \code{std_logic} & Den avstoppade bit som accepterades. \\
  11 & \code{real_bit_valid} & out & \code{std_logic} & \code{'1'} när en riktig bit accepterades,
    \code{'0'} för en kastad stoppbit. \\
\end{longbooktable}

Sändarens spegelbild. Efter \code{MAX_RUN} lika bitar i rad \emph{måste} nästa bit vara den motsatta;
är den det kastas den och räkningen börjar om, är den det inte är det ett \code{stuff_error}.

\subsection*{\code{can_controller} (\chapref{c11:ch}, \ref{c16:ch}, \ref{c17:ch})}
Toppnivån för CAN-halvan, och det enda block som känner till ramformatet. Femton portar, i exakt den
här ordningen, alla ingångar först:

\begin{textdrawing}
 1 clock     4 tx_id    7 rx_bus    10 bus_en   13 rx_data
 2 reset_n   5 tx_dlc   8 tx_done   11 rx_id    14 rx_valid
 3 tx_req    6 tx_data  9 tx_bus    12 rx_dlc   15 error
\end{textdrawing}

\textbf{Ramsekvensen} en sändning går igenom: SOF (1 bit), arbitreringsfältet (11 bitars
identifierare + RTR), kontrollfältet (IDE, r0, DLC = 6 bitar), datafältet (0--64 bitar enligt DLC),
CRC (15 bitar + avgränsare), ACK (lucka + avgränsare) och EOF (7 bitar). Varje fält får ett eget
laddningstillstånd före sitt skifttillstånd.

\textbf{Bussen} modelleras som öppen dränering: \code{bus_en} avgör om noden driver alls,
\code{tx_bus} vad den driver. Bussvärdet är \code{'1'} bara om ingen nod drar ned det. Under
ACK-luckan släpper den sändande noden bussen, och en mottagare med godkänd CRC drar ned den.
\textbf{Sändaren läser aldrig tillbaka ACK-luckan}; en ram som ingen kvitterar fullbordas precis som
en kvitterad, och \code{tx_done} pulsar ändå.

\textbf{Arbitrering:} under arbitreringsfältet jämförs varje sänd bit med det som faktiskt ligger på
bussen. Skickar noden en recessiv bit och läser tillbaka en dominant har den förlorat: den sätter
\code{error}, pulsar \code{tx_done}, lämnar sändarrollen och går till tomgång. Den \textbf{tar inte}
emot ramen som vann.

Att \code{tx_done} pulsar även här är avsiktligt: pulsen betyder ”sändningsförsöket är över och
porten är din igen”, inte ”ramen kom fram”. \code{error} säger vilketdera. Utan den pulsen har en
anropare ingen flank att vänta på, och registerbanken får aldrig tillbaka sin TX-klar-bit.

\textbf{\code{error}} går hög av fyra orsaker: förlorad arbitrering, \code{stuff_error}, misslyckad
CRC-kontroll, och en mottagen fjärrram. Den ligger kvar tills nästa sändning eller mottagning
startar. \textbf{\code{tx_done}} och \textbf{\code{rx_valid}} är pulser på exakt en klockcykel, vilket
är precis det som gör registerbanken nödvändig.

\subsection{Kända begränsningar som drivrutinen märker av}
\label{app:project:limits}

Förenklingarna mot riktig CAN listas i sin helhet i \secref{c19:sec:gaps}, och från protokollhållet i
\canref{7}. Tre av dem syns ända ut i registerkartan och måste därför stå här också, eftersom
parallellklassen skriver kod mot dem \textemdash{} det är samma tre som \canbook{} lyfter fram, och
den boken läser båda klasserna:

\begin{itemize}
  \item \textbf{En förlorad arbitrering går inte att skilja från en lyckad sändning på STATUS
    bit 0.} Båda ger tillbaka biten, eftersom kontrollern pulsar \code{tx_done} också när den
    avbryter. Det är avsiktligt: biten betyder ”porten är fri”, inte ”ramen kom fram”. En drivrutin
    som vill veta utfallet läser \code{ERROR_FLAGS} efter att bit 0 kommit tillbaka.

    Det här var en gång en riktig bugg, och det är värt att veta varför den inte är det längre.
    Sattes bit 0 enbart av \code{tx_done}, och pulsade \code{tx_done} bara vid EOF, blev biten
    liggande låg efter varje förlorad arbitrering och noden kunde inte sända igen förrän den
    nollställdes på resetknappen \textemdash{} på en tvånodsbuss alltså i normal drift. Tre saker gör
    att den inte kan uppstå igen, och de är medvetet överlappande: kontrollern pulsar \code{tx_done}
    även vid avbrott (\chapref{c17:ch}), banken sätter bit 0 också på en stigande flank av
    \code{error} medan biten är låg (\chapref{c18:ch}), och \code{TX_ABORT} finns som nödutgång.
    \code{can_controller_tb} fall 2 och \code{register_bank_tb} fall 10--11 kontrollerar de tre.
  \item \textbf{\code{error} är en bit för fyra orsaker.} Drivrutinen kan se \emph{att} något gick
    fel, aldrig \emph{vad}. Omsändningslogik som skiljer på förlorad arbitrering och trasig CRC går
    inte att skriva mot det här gränssnittet.
  \item \textbf{Ingen buffring på mottagarsidan.} Kontrollern håller en ram. Kommer nästa innan den
    förra lästs ut skrivs den över. Registerbanken vidarebefordrar begränsningen i stället för att
    dölja den; se bilaga~\ref{app:protocol}.
\end{itemize}

% ------------------------------------------------------------------------------------------
\section{Modulerna i \code{bridge/}}
\label{app:project:bridge}

\subsection*{\code{register_bank} (\chapref{c18:ch})}
Sexton portar, i den här ordningen:

\begin{textdrawing}
 1 clock        5 reg_wdata    9 rx_data     13 tx_id
 2 reset_s2_n   6 tx_done     10 rx_valid    14 tx_dlc
 3 reg_index    7 rx_id       11 error       15 tx_data
 4 reg_write    8 rx_dlc      12 reg_rdata   16 tx_req
\end{textdrawing}

Översätter kontrollerns encykelspulser till klibbiga, pollbara nivåer, och registerskrivningar till
händelser. Semantiken i sin helhet står i bilaga~\ref{app:protocol}; den är kontraktet mot
parallellklassen och får inte ändras ensidigt. \code{reg_rdata} är \textbf{kombinatorisk}: att låsa
ett läst värde är bryggans ansvar, inte bankens.

\subsection*{\code{spi_reg_bridge} (\chapref{c19:ch})}
Tio portar, i den här ordningen: \code{clock}, \code{reset_s2_n}, \code{ss_active}, \code{rx_data},
\code{rx_valid}, \code{reg_rdata}, \code{tx_data}, \code{reg_addr}, \code{reg_wdata},
\code{reg_write}.

Implementerar transaktionen ur bilaga~\ref{app:protocol}: en kommandobyte plus fyra databytes, alltid
exakt fem. Vid läsning \textbf{låses registervärdet en gång} i slutet av kommandobyten; vid skrivning
verkställs ingenting förrän femte byten är hel. Går \code{ss_active} låg dessförinnan avbryts
transaktionen helt, utan sidoeffekter. \code{reg_addr} \textbf{håller} det avkodade indexet tills
nästa kommandobyte avkodar ett nytt.

\subsection*{\code{can_spi_node} (\chapref{c19:ch})}
Nio portar, i den här ordningen:

\begin{textdrawing}
 1 clock      4 mosi     7 rx_bus
 2 reset_n    5 ss       8 tx_bus
 3 sclk       6 miso     9 bus_en
\end{textdrawing}

Rent strukturell toppnivå. Instansierar \code{spi_slave}, \code{spi_reg_bridge},
\code{register_bank} och \code{can_controller}, plus en \code{meta_prev} för reset. Ingen funktionell
logik här.

Ordningen är SPI:s fyra ledningar som en grupp, \code{sclk}/\code{mosi}/\code{ss}/\code{miso}, och
därefter CAN-bussens tre. \code{miso} ligger alltså bland SPI-portarna och inte hos de andra
utgångarna, så att de fyra pinnarna står i samma ordning som i kopplingstabellen. Det spelar roll:
\code{can_spi_node_tb} binder positionellt som allt annat, och \code{miso} och \code{rx_bus} har
samma typ, så en förväxling av dem passerar analysen utan ett ord och ger en nod som aldrig svarar.

% ------------------------------------------------------------------------------------------
\section{Testplan}
\label{app:project:tests}

Gruppen skriver inga testbänkar i det här projektet. Nio är utdelade, och de är kontraktet: sju
modultestbänkar och två systemtestbänkar.

\begin{booktable}{lL}
  \thead{Testbänk} & \thead{Kontrollerar} \\
  \midrule
  \code{meta_prev_tb} & Två vippors fördröjning, för både en bit och en vektor. \\
  \code{bit_timer_tb} & \code{sample} och \code{bit_done} vid rätt tick; \code{resync};
    \code{enable}. \\
  \code{crc15_tb} & CRC-värden för kända sekvenser; \code{valid}; \code{clear}. \\
  \code{tx_shift_reg_tb} & MSB-först skiftning, stoppning efter fem lika, \code{done}-tajmingen. \\
  \code{rx_shift_reg_tb} & Deserialisering, avstoppning, \code{stuff_error}, \code{real_bit}. \\
  \code{register_bank_tb} & Elva fall, från resettillståndet till \code{TX_ABORT} och de två vägarna
    tillbaka till TX klar. \\
  \code{spi_reg_bridge_tb} & Fyra fall: en skrivning, en läsning byte för byte, ett avbrott på
    \code{ss_active}, och en kommandobyte med reserverad bit satt. \\
\end{booktable}

\code{can_controller_tb} sätter två noder på en gemensam wired-AND-buss och kontrollerar att den ena
tar emot den andras ram, och att arbitrering löses upp rätt när båda sänder samtidigt. Den hoppas
över av bygget tills kontrollern har en mottagarväg (bilaga~\ref{app:sim}).

\code{can_spi_node_tb} gör samma sak ett lager ut: två kompletta noder på samma buss, men varje nod
driven av en modellerad SPI-master på sina fyra pinnar. Den spelar upp protokollspecifikationens
genomarbetade exempel genom hela kedjan och följer en ram från en SPI-skrivning på den ena noden till
en SPI-läsning på den andra. Sju fall: resettillståndet, det genomarbetade exemplet, ogiltiga
kommandobytes och reserverade index, \code{TX_ABORT}, avbrottsregeln, extra bytes under samma
\code{SS}, och till sist sändningen och mottagningen.

Det är den \textbf{enda} testbänk i kursen som driver \code{sclk}, \code{mosi} och \code{ss} som
riktiga vågformer. \code{spi_reg_bridge_tb} matar bryggan färdiga bytes och rör aldrig en pinne, så
allt mellan ledningen och byten \textemdash{} synkroniserarna, flankdetekteringen på SCK,
\code{ss_active}-ramningen \textemdash{} provas här eller ingenstans.

\subsection{Egen verifiering}
\label{app:project:ownverify}

\code{can_spi_node_tb} är snäll mot tidskraven med flit: den håller 200 ns tyst kring varje
\code{SS}-flank, där protokollet kräver minst 60 ns. Det är rätt för en testbänk som ska skilja
logikfel från tidsfel, men det betyder också att \textbf{ingen} utdelad testbänk pressar
specifikationens tidsminima.

Det är den verifiering projektet ber gruppen göra själv: kopiera \code{can_spi_node_tb}, krymp den
tysta perioden ner mot 60 ns, och ta reda på vid vilket värde noden slutar svara och varför. Svaret
ligger i \code{spi_slave}s synkroniserare, och det är det enda stället där kursen låter dig
\emph{mäta} en marginal i stället för att läsa den. Övning~\ref{c19:ex:node}(b) ställer frågan i sin
helhet.

\subsection{Checklista före inlämning}
\label{app:project:checklist}

\begin{itemize}
  \item Alla nio utdelade testbänkar passerar, \code{can_controller_tb} och \code{can_spi_node_tb}
    inräknade. \textbf{Kontrollera att de verkligen \emph{kördes} och inte hoppades över}, se rutan
    nedan.
  \item \code{make build-project} går grönt från ett rent utcheckat repo.
  \item Portordningen stämmer mot tabellerna i \secref{app:project:controller} och
    \ref{app:project:bridge}.
  \item Inga värden hårdkodade förbi \code{can_def}.
  \item Varje entitets gränssnitt är kommenterat.
  \item \file{CONTRIBUTORS.md} är aktuell.
  \item \code{main} innehåller inga commits som inte gått via en granskad PR.
\end{itemize}

\begin{warning}
\textbf{Överhoppad är inte godkänd.} \code{can_controller_tb} grindas på om
\file{can_controller.vhd} innehåller \chapref{c17:ch}:s CRC-grindning, letad upp med en textsökning
på de interna signalnamn kapitlet använder. Skriver ni en korrekt mottagarväg med andra namn
rapporteras testbänken som \emph{skipped}, och eftersom överhoppad aldrig är underkänd ser
checklistan uppfylld ut medan kursens pass/fail-testbänk aldrig kördes. Kör därför

\begin{shell}
CI_BUILD_ALL=1 make build-project
\end{shell}

minst en gång före inlämning, och lämna in det utfallet.
\end{warning}

% ------------------------------------------------------------------------------------------
\section{Examination och betyg}
\label{app:project:grading}

Kursen examineras genom \textbf{ett grupprojekt} och \textbf{två praktiska tentamina}. Totalt 100
poäng.

\begin{booktable}{lllr}
  \thead{Moment} & \thead{Var} & \thead{Form} & \thead{Poäng} \\
  \midrule
  Praktisk tentamen 1, sekvensnät & \chapref{c9:ch} & Individuell, 3 h & 25 \\
  Grupprojekt, en CAN-nod i VHDL & \chapref{c10:ch}--\ref{c19:ch} & Grupp om 4--5 & 50 \\
  Praktisk tentamen 2, CAN-moduler & \chapref{c20:ch} & Individuell, 3 h & 25 \\
\end{booktable}

\begin{booktable}{lL}
  \thead{Betyg} & \thead{Krav} \\
  \midrule
  \textbf{IG} & Under 50 poäng, eller ett underkänt projekt. \\
  \textbf{G} & Minst 50 poäng totalt \textbf{och} godkänt projekt. \\
  \textbf{VG} & Minst 80 poäng totalt \textbf{och} väl godkänt projekt. \\
\end{booktable}

Projektet är alltså en spärr i båda riktningarna. Det står för halva kursen, det är där det mesta av
arbetet ligger, och ett underkänt projekt går inte att kompensera med två starka tentor.

\subsection{Projektets 50 poäng}
\label{app:project:points}

\begin{booktable}{LrL}
  \thead{Del} & \thead{Poäng} & \thead{Vad som bedöms} \\
  \midrule
  Designen & 30 & Modulerna fungerar enligt specifikation och testbänkarna passerar. \\
  Arbetssättet & 12 & Git-historik, PR:er, granskningarnas kvalitet, seminariet,
    \file{CONTRIBUTORS.md}. \\
  Redovisningen & 8 & Slutredovisningen: att gruppen kan förklara vad den byggt och varför. \\
\end{booktable}

Betyget sätts på gruppen. Är insatsen påtagligt ojämn kan enskilda betyg justeras; det är den enda
anledningen till att \file{CONTRIBUTORS.md} ska hållas aktuell löpande och inte skrivas kvällen före.

\textbf{Godkänt projekt} kräver en \textbf{komplett nod}: alla nio utdelade testbänkar ska passera,
alltså \code{CI_BUILD_ALL=1 make build-project} grönt rakt igenom. Kontrollern ensam räcker inte, och
skälet ligger efter kursen: noden följer med till CAN-labben i nästa kurs i inbyggda system, där en
drivrutin ska nå den över SPI. En kontroller utan registerbank och brygga går inte att nå från någon
drivrutin alls, och en nod vars delar fungerar var för sig men inte ihop går inte heller att nå.

\textbf{Väl godkänt} kräver dessutom den egna verifieringen i \secref{app:project:ownverify}, att
koden är genomgående dokumenterad och stilistiskt konsekvent, och att granskningskulturen syns i
repots historik.

Bring-up på riktig hårdvara bedöms \textbf{inte}. Ett trasigt kit eller en trasig koppling får inte
kosta poäng. Allt som bedöms går att köra i GHDL på en vanlig laptop.

\subsection{De praktiska tentorna}
\label{app:project:exams}

Båda är individuella, tre timmar, vid datorn med GHDL. Formen är densamma i båda: fristående
uppgifter, oberoende av varandra och ordnade ungefär efter stigande svårighet.

De flesta uppgifterna är \textbf{skrivuppgifter}: en utdelad, självkontrollerande testbänk och en
beskrivning av modulen den driver, med portordning och beteende. \Chapref{c20:ch} har därutöver
\textbf{läsuppgifter}: ett tidsdiagram eller en SPI-transaktion att tolka på papper. De har ingen
testbänk, av precis det skäl som gör dem värda att ställa \textemdash{} att läsa en vågform är en
färdighet i sig, och den examineras inte av att skriva VHDL. De rättas mot ett facit, och delvis rätt
svar ger delpoäng på samma sätt som delvis korrekt logik gör.

Delvis korrekt logik ger delpoäng: lämna aldrig in en tom fil för att modulen inte blev klar, och
samma sak gäller ett halvt uttolkat tidsdiagram. Rättningen läser koden, inte bara testbänkens
utfall; en modul som passerar genom att göra något annat än det som efterfrågades ger inte full poäng.

\Chapref{c9:ch} omfattar \chapref{c1:ch} till \ref{c8:ch}, med tyngdpunkt på sekvensnät: register,
räknare, skiftregister, timers, synkroniserare och enkla Mooremaskiner. Ett övningsprov delas ut i
förväg. \Chapref{c20:ch} omfattar projektets moduler: bittajming, CRC, bitstoppning åt båda hållen,
ramsekvensering, registersemantik och SPI-transaktionen.

\subsection{Hjälpmedel}
\label{app:project:aids}

Det här avsnittet är den \textbf{auktoritativa} listan. Står något annat någon annanstans gäller det
som står här.

\textbf{Tillåtet vid båda tentorna:}
\begin{itemize}
  \item Allt kursmaterial: den här boken, föreläsningarnas README och samtliga appendix. Vid
    \chapref{c20:ch} hör registerkartan och protokollspecifikationen (bilaga~\ref{app:protocol})
    förstås dit; vid \chapref{c9:ch} finns de också tillgängliga, men examineras inte.
  \item All VHDL du själv skrivit under kursen, och vid \chapref{c20:ch} även gruppens projektkod.
  \item Egen dator med GHDL, och CircuitVerse. Båda tentorna: att rita en krets innan man skriver
    den är ofta snabbare.
\end{itemize}

\textbf{Inte tillåtet:}
\begin{itemize}
  \item \textbf{Kommunikation med andra}, i någon form.
  \item \textbf{Verktyg som genererar eller förklarar VHDL åt dig}: språkmodeller, kodassistenter och
    liknande, oavsett om de skriver koden eller bara resonerar om den. Gränsen går inte vid vem som
    trycker på tangenterna utan vid vems förståelse som prövas.
  \item \textbf{Sökning på nätet efter lösningar.} CircuitVerse körs i en webbläsare, så datorn är
    uppkopplad, och det förutsätter att du håller dig till kursmaterialet av egen kraft. Att slå upp
    GHDL:s eller VHDL:s dokumentation är tillåtet; att söka efter en färdig modul är det inte.
\end{itemize}

Är du osäker på om något är tillåtet: fråga före tentan, inte under.

Att gruppens projektkod är tillåten vid \chapref{c20:ch} är ett medvetet val. Ni byggde den.
Uppgifterna är i stället skrivna som varianter och reparationer, så att de skiljer sig från
projektets moduler på just den punkt som kräver att man förstått dem.

\subsection{Omexamination och frånvaro}
\label{app:project:reexam}

Båda tentorna går att skriva om. Projektet lämnas in en gång, med möjlighet till komplettering efter
besked från handledaren: en komplettering rättas mot godkänt, inte mot väl godkänt.

Kodgranskningsseminariet i \chapref{c15:ch} och slutredovisningen i \chapref{c19:ch} är projektets två
obligatoriska hållpunkter. Missad närvaro tas igen efter överenskommelse med handledaren. Övriga pass
har ingen närvaroplikt; grupprojektet har däremot en gruppdynamik som inte tål att någon uteblir,
vilket är ett annat och starkare skäl att komma.

% ------------------------------------------------------------------------------------------
\section{Praktiska tips}
\label{app:project:tips}

\begin{itemize}
  \item Rita tidsdiagrammet innan ni skriver skiftregistren. De flesta fel där är tajmingfel på en
    cykel, inte logikfel, och de syns på papper innan de syns i simulatorn.
  \item Börja med \code{bit_timer} och \code{crc15} om ni är osäkra på var ni ska börja. De är de
    enklaste, de beror inte på något, och de går att bli helt klar med.
  \item Läs de utdelade testbänkarna. De säger mer exakt vad varje modul lovar än någon prosa gör,
    och de är skrivna för att läsas.
  \item Testa stoppningen noggrant, åt båda hållen. Det är projektets vanligaste felkälla.
  \item Håll \code{can_spi_node} strukturell. Så fort det smyger sig in logik i toppnivån blir den
    omöjlig att testa isolerat.
  \item Prata med parallellklassen tidigt, och gör det om registerkartan snarare än om er kod.
\end{itemize}

\begin{aside}
\textbf{Gemensam dokumentation.} Två dokument delas med mjukvarugruppen och är kontraktet mellan
klasserna: registerkartan (vad registren betyder) och protokollspecifikationen (hur de nås), båda
återgivna i bilaga~\ref{app:protocol}. Ändras något i dem måste båda klasserna veta om det innan
ändringen mergas. Dokumenten uppdateras \emph{före} koden, aldrig efter.
\end{aside}
